\documentclass[11pt]{article}

% Standard, arXiv-safe packages only (compiles with a plain pdflatex + one bibtex-free run, or tectonic).
\usepackage[utf8]{inputenc}
\usepackage[T1]{fontenc}
\usepackage{amsmath,amssymb}
\usepackage{booktabs}
\usepackage{array}
\usepackage{graphicx}
\usepackage[margin=1in]{geometry}
\usepackage{microtype}
\usepackage[hidelinks]{hyperref}
\usepackage{xcolor}

% NOTE TO AUTHOR (Daniel): this is a findings-paper draft. The three previously-flagged citations
% (Control Tax arXiv:2506.05296; DeepMind Frontier Safety Framework v3.0, 22 Sep 2025; METR MALT,
% 14 Oct 2025) were verified against primary sources on 2026-07-14. Before sending anywhere, still confirm
% the author/affiliation block and re-read the Limitations section. Every quantitative claim below
% regenerates from a committed runner, manifest, or validated receipt in the Telos repository; the exact
% experiment is named in the text (for example, iter195).

\title{Reward Hacks That Pass the Test Suite:\\
Constructing and Detecting SWE-bench Patches\\
That Are Certified Yet Wrong}

\author{Daniel Wahnich\\ \texttt{Aweb} \\ \small Telos research program}

\date{\today}

\begin{document}
\maketitle

\begin{abstract}
A coding agent is usually judged by whether a test suite passes. Passing the tests is not the same as
solving the task: an agent can make the graded tests pass while leaving the code wrong on inputs the tests
never check. In reinforcement-learning terms, the test suite is a proxy reward, and passing it without
solving the task is reward hacking. We ask whether SWE-bench Verified, the standard benchmark for coding
agents, certifies patches that are actually wrong, and if so, what detects them.

We report three results. First, a self-correction: an earlier benchmark we built did not contain reward
hacks at all. Its generator only kept a patch when the patch broke the project's existing test suite, so
the official grader marks all forty rows unresolved, and simply running the existing tests catches all
forty at no model cost. Second, we construct the real class. Using a frontier model as an adversary, we
build ten patches that the official SWE-bench harness marks resolved---each passes every graded test---yet
each produces different output from the correct fix on an input the tests do not cover, verified by
execution in the project's own container. We then extend the set to twenty-two such patches across eight
repositories, so the class is not specific to a few projects. Third, we compare two detectors at decision
time on the first ten hacks, with the correct fix withheld from both. A panel of two frontier
language-model judges, reading only the patch, flags seven of ten hacks but also flags correct patches. A
gold-free oracle, which proposes an executable correctness check from the public task description and runs
it, flags four of ten and never flags a correct patch. The two detectors catch different hacks; together
they flag eight of ten. An ensemble of a judge and a precise execution check therefore beats either
detector alone.

These are small, deliberately constructed pilots (a benchmark of twenty-two hacks across eight
repositories, with the detector comparison run on the first ten, one adversary model), not a survey of how
often agents hack in practice. Every number regenerates from committed
artifacts.
\end{abstract}

\section{Introduction}

An autonomous coding agent is given a bug report and asked to produce a patch. The usual measure of success
is a test suite: if the tests pass, the task is considered done. This measure is convenient, and it is also
gameable. The agent can edit or weaken the tests, or it can write source code that handles only the exact
input the visible test checks rather than solving the general problem. When the test suite is a proxy for
``the task is solved,'' making the tests pass without solving the task is \emph{reward hacking}, the general
failure in which an optimizer satisfies a proxy reward while missing the objective the reward stands
in for~\cite{amodei2016concrete,krakovna2020specification,skalse2022defining}.

This paper asks a specific version of that question. SWE-bench Verified is the standard benchmark for
coding agents~\cite{jimenez2024swebench}. Each of its instances comes with two sets of tests: a
\texttt{FAIL\_TO\_PASS} set that a correct fix must newly pass, and a \texttt{PASS\_TO\_PASS} set of
existing tests that a correct fix must keep passing. The official harness marks an instance \emph{resolved}
when a patch passes both sets. Our question is: can a patch pass both sets---be certified resolved---and
still be wrong? If it can, then no test-based check can catch it, because those tests are exactly what
``resolved'' means.

The question matters beyond one benchmark. Google DeepMind's Frontier Safety Framework describes a model
that circumvents oversight ``even when relevant model outputs \ldots{} are being monitored,'' and lists the
mitigation for that case as future work~\cite{deepmind2025fsf}. A patch that passes every graded test while
being wrong is a concrete, measurable instance of that case.

We make three contributions.
\begin{enumerate}
  \item \textbf{A self-correction (Section~\ref{sec:correction}).} An earlier benchmark we released did
  not contain reward hacks. We show why from committed evidence, and we report the free baseline that had
  been missing.
  \item \textbf{Certified-yet-wrong patches (Section~\ref{sec:construct}).} We construct and release a
  benchmark of twenty-two patches across eight repositories that the official SWE-bench harness marks
  resolved and that execution shows are wrong, each with a retained witness.
  \item \textbf{A two-detector comparison (Section~\ref{sec:detect}).} We measure a language-model judge
  panel and a gold-free execution oracle on the ten hacks, with the correct fix withheld from both, and
  report a recall, precision, and coverage tradeoff, plus the finding that the two detectors are
  complementary.
\end{enumerate}

\section{Background}

\paragraph{Reward hacking.} Reward hacking, also called specification
gaming~\cite{krakovna2020specification}, is when an agent scores well on a proxy reward while failing the
real goal. Skalse et al.~\cite{skalse2022defining} give a formal definition. We use the term in its plain
sense: a patch that scores well on the grader (it is certified resolved) while being wrong. We note one
name collision to avoid confusion: RewardBench~\cite{lambert2024rewardbench} evaluates \emph{reward
models}, which rank preferences, and is unrelated to the reward hacking studied here.

\paragraph{SWE-bench Verified.} SWE-bench~\cite{jimenez2024swebench} asks a model to resolve a real GitHub
issue in a real Python repository. SWE-bench Verified is a human-filtered subset. Each instance pins a
repository commit and a container image, and defines the \texttt{FAIL\_TO\_PASS} and \texttt{PASS\_TO\_PASS}
test sets described above. The ``gold patch'' is the accepted human fix. A known weakness of the format is
that the graded tests can be incomplete: they may not exercise every behavior a correct fix must have. Our
constructed hacks live in exactly that gap.

\paragraph{Language-model judges.} A language model can act as a monitor that reads a patch and decides
whether it is correct~\cite{zheng2023judging}. Judges are not uniformly weak. On backdoored programming
solutions, a monitor can reach an area-under-the-curve of 0.92 to 0.96~\cite{terekhov2025controltax}. Their
accuracy drops on subtle, semantic errors in real repositories. Our judge result---seven of ten---is
consistent with judges having real but imperfect signal, not with judges being useless.

\paragraph{Gold-free and execution-based checks.} A different family of checks does not read the patch and
judge it, but instead constructs an input, runs the code, and checks the output. When such a check does not
use the correct reference fix at decision time, we call it \emph{gold-free}. This family includes
property-based testing~\cite{claessen2000quickcheck} and metamorphic testing~\cite{chen2018metamorphic}.
Our second detector is a gold-free execution check.

\paragraph{Scale of existing corpora.} Public datasets of reward hacking are far larger than ours: METR's
MALT contains on the order of ten thousand transcripts~\cite{metr2025malt}. Ours is ten hacks. We make no
claim of scale. What is different here is the unit and the verification: we label individual patches (not
whole trajectories), and every label is verified by execution. Like the larger corpora, ours is elicited:
a model was asked to produce the hacks. We therefore claim diff-level, execution-verified construction, not
naturally occurring behavior.

\section{An Earlier Benchmark Measured the Wrong Thing}
\label{sec:correction}

We begin with a correction, because it defines the target for the rest of the paper.

An earlier stage of this project released a forty-row artifact and described its rows as reward hacks. The
generator that produced those rows kept a candidate only when the candidate made at least one
\texttt{PASS\_TO\_PASS} test fail. \texttt{PASS\_TO\_PASS} is the project's own existing test suite. So
every row broke the existing suite by construction.

We re-derived the official harness verdict for all forty rows from the committed evaluation logs. All forty
are marked unresolved, and all forty have at least one failing \texttt{PASS\_TO\_PASS} test. The
consequence is direct. Running the existing test suite---equivalently, reading the official harness
verdict---detects all forty rows at zero model cost. A three-model language-model judge panel, the detector
we had been studying, detected seventeen of forty at a metered cost of \$13.59. That free baseline had
never been reported.

The rows are therefore not reward hacks. A reward hack scores well on the reward while being wrong; these
rows fail the reward. They are wrong patches that break the existing tests, which any test run catches for
free. The earlier judge figure of seventeen of forty is arithmetically correct and we do not retract it;
what changes is what it measures. It measures how often a judge, reading a diff, spots a wrong patch that
the existing tests already reject. That is a fact about code review, not about reward hacking, and it does
not motivate a new verifier.

The real target is the inverse: a patch that \emph{passes} the whole suite and is still wrong. The official
harness scores such a patch as resolved, so running the tests cannot catch it. The rest of the paper
constructs that class and asks what does catch it.

\section{Constructing Certified-Yet-Wrong Patches}
\label{sec:construct}

\paragraph{Definition.} A \emph{certified-resolved reward hack} is a patch that the official SWE-bench
harness marks resolved---it passes every \texttt{FAIL\_TO\_PASS} and every \texttt{PASS\_TO\_PASS}
test---and that produces different output from the gold patch on some input the graded tests do not cover.

\paragraph{Why the shipped tests cannot witness these.} Before constructing the class, we measured whether
the project's own tests could ever reveal it. We took sixteen candidate patches, confirmed by execution
that all sixteen are certified resolved, and then ran, for each one, every test in the module that contains
its graded tests. For the Django instances, the graded \texttt{PASS\_TO\_PASS} set is the entire test
module: there are no additional tests to hold out. Across all sixteen instances, the module tests beyond
the graded set number 193, and not one of the 193 distinguishes a certified patch from the gold patch. A
certified patch passes every shipped test in its module, graded or not. To witness that such a patch is
wrong, a check must run an input that no shipped test runs.

\paragraph{Construction.} We construct the class in three steps, all executed in each instance's pinned
container.
\begin{enumerate}
  \item \textbf{Elicit a variant.} We take the gold patch and ask a frontier model to rewrite its largest
  added block into a variant that still passes the graded tests but is subtly wrong on other inputs. The
  model is explicitly asked to produce a wrong variant; the class is elicited, not observed in the wild.
  \item \textbf{Certify by execution.} We run the instance's official SWE-bench evaluation script in its
  pinned container---\texttt{runtests.py} for Django, \texttt{pytest} for the others---and parse the
  result with the official SWE-bench log parsers. A variant is kept only if it passes every
  \texttt{FAIL\_TO\_PASS} and \texttt{PASS\_TO\_PASS} test.
  \item \textbf{Witness the wrongness.} To label a certified variant as wrong, we use the gold patch as
  ground truth---but only here, at labelling time, never at detection time. A model constructs a small
  scenario that calls the changed function on an input designed to reach the branch where the variant and
  the gold fix could differ. We run that scenario under the gold patch and under the variant, in the same
  container. If the two produce different output, and the gold run is clean, the variant is a witnessed
  reward hack.
\end{enumerate}

\paragraph{Result.} From sixteen certified candidates, this procedure yields ten witnessed
certified-resolved reward hacks across three repositories (Django, astropy, matplotlib). Two further
candidates turned out to be correct rewrites---the variant matched the gold patch on every input we
tried---and we report them as such rather than discard them. The remaining candidates produced no clean
witness. Table~\ref{tab:hacks} lists these ten hacks, which are the set used for the detector comparison in
Section~\ref{sec:detect}.

\paragraph{The class is not specific to three repositories.} We then ran the same pipeline on forty-two
fresh instances spanning all twelve repositories in SWE-bench Verified. It produced twenty-three variants;
twenty were certified resolved by execution, and twelve were also witnessed wrong, adding hacks in five new
repositories: requests, xarray, pylint, pytest, and sympy. Combined with the first ten, the released
benchmark is twenty-two execution-verified certified-resolved reward hacks across eight repositories. The
sharpest scale example is a Django hack whose variant passes 1{,}870 graded tests and still returns
\texttt{'accepted'} where the correct fix returns \texttt{'rejected'}. The detector comparison below was
run on the first ten hacks only, before this expansion.

\paragraph{The class also occurs unprompted.} Every hack described so far is elicited: a model was told to
produce a wrong variant. We therefore ran one further check with no such instruction. We asked a frontier
model to fix thirty-nine issues with a plain prompt---solve this bug---that never mentions tests, and
measured how often its patch is certified yet wrong. Of thirty-seven patches, fifteen were certified
resolved by the official harness; eight matched the correct fix, and seven produced different output. To
decide whether a difference is a genuine bug or a valid alternative fix, two judges saw the two outputs
unlabeled, neither told which came from the reference; we counted a patch as wrong only when both judges
named that patch, and only that patch, as incorrect. One patch met this bar: a Sphinx fix that passes all
thirty-three graded tests yet renders a compound-key keyboard input as a malformed empty element. This is a
single case on a localized-solve sample, an existence result and not a frequency: a model, not told to game
the tests, nonetheless produced a patch that passes them while being wrong. Two further cases were
ambiguous (a judge doubted the reference fix as well) and are not counted.

\begin{table}[t]
\centering
\small
\caption{The ten certified-resolved reward hacks. Each patch passes every graded test (the official
harness marks it resolved), yet produces different output from the gold fix on the shown input.}
\label{tab:hacks}
\begin{tabular}{@{}llll@{}}
\toprule
Instance & Function & Gold output & Variant output \\
\midrule
django-11179 & \texttt{Collector.delete}        & \texttt{None} & \texttt{'present'} \\
django-11119 & \texttt{Engine.render\_to\_string} & \texttt{<tag>} & \texttt{\&lt;tag\&gt;} \\
django-11163 & \texttt{model\_to\_dict}          & \texttt{\{\}} & \texttt{\{'id': None, ...\}} \\
django-11141 & \texttt{MigrationLoader.load\_disk} & \texttt{(True, False)} & \texttt{(False, True)} \\
django-11211 & \texttt{UUIDField.deconstruct}    & \texttt{UUID('550e...')} & \texttt{'550e...' (str)} \\
django-11433 & \texttt{construct\_instance}      & \texttt{0} & \texttt{7} \\
astropy-7166  & \texttt{is\_public\_member}       & \texttt{'...documentation'} & \texttt{None} \\
astropy-14096 & \texttt{SkyCoord.\_\_getattr\_\_}  & \texttt{'intercepted'} & \texttt{AttributeError} \\
matplotlib-23412 & \texttt{Patch.draw}           & \texttt{((2.5, ...),)} & \texttt{((0.0, ...),)} \\
matplotlib-24627 & \texttt{\_AxesBase.\_\_clear}  & \texttt{(True, True)} & \texttt{(False, False)} \\
\bottomrule
\end{tabular}
\end{table}

\paragraph{A worked example.} The clearest hack is \texttt{django-11179}. When a database object is deleted,
the correct fix clears the object's primary key to \texttt{None} in memory. The gold patch does exactly
that. The variant clears the key only when the key is literally named \texttt{id}:
\begin{center}
\texttt{setattr(instance, pk.attname, None if pk.name == 'id' else instance.pk)}
\end{center}
Every graded test for this instance uses a model whose primary key is named \texttt{id}, so the variant
passes them and the harness certifies it. But a model whose primary key has any other name keeps a stale
key after deletion. Running the scenario, the gold patch prints \texttt{None} and the variant prints the
stale value.

\paragraph{Execution beat our own judgment.} One case is worth stating against ourselves. We hand-analyzed
\texttt{django-11119} and judged its variant a correct rewrite: the changed expression looked equivalent
for boolean inputs. Execution disagreed. With auto-escaping turned off, the gold patch renders text
unescaped and the variant escapes it, so the two differ. A reasoning pass called the patch correct; running
the code showed it was a hack. This is the paper's own thesis applied to the authors: for this class, run
the code rather than trust a verdict, including one's own.

\paragraph{Claim boundary.} The ten hacks are elicited and constructed, not sampled from real agent runs,
so we make no claim about how often agents hack in practice. There are twenty-two hacks across eight
repositories, and the detector comparison uses the first ten.
``Wrong'' means the variant differs from the gold reference fix on a constructed input. The two correct
rewrites are reported, not hidden.

\section{Detecting Certified-Yet-Wrong Patches}
\label{sec:detect}

We now ask what catches the ten hacks. Every detector sees only the candidate patch, the public task
description, and the visible test. No detector sees the gold patch, the label, the hidden tests, or the
witness scenario from Section~\ref{sec:construct}.

\paragraph{Controls.} Recall without a false-positive rate is not a result. Alongside the ten hacks we
evaluate two kinds of control: the ten gold patches for the same instances, which are correct, and the two
correct rewrites from Section~\ref{sec:construct}, which are a harder control because they look unusual
while being correct.

\paragraph{Detector A: a judge panel.} Two frontier language models act as reviewers: they read the patch,
the task, and the visible test, and answer whether the patch is a correct general fix or subtly wrong. We
had intended a three-model panel; the third model was unreachable in our environment, so we report a
two-model panel and say so. We count a hack as caught if at least one judge flags it. On a single run, the
panel flags seven of the ten hacks, one of the ten gold patches, and one of the two correct rewrites.
Language-model judge outputs vary from run to run, so this is one sample, not a stable rate.

Two measurement problems arose and we fixed both before reporting. Our first automatic check for leaked
information matched a shared 200-character diff header and reported leakage on all ten hacks; this was
wrong, because a variant shares its diff header with the gold patch by construction. The corrected check
counts only gold-distinctive lines that appear in the candidate, which is two of ten, and a shared correct
line makes a hack look \emph{more} correct, so it biases the judge toward missing the hack, not toward
catching it. Separately, one judge's short token budget caused it to emit unparseable output on most rows;
counting those as ``correct'' would have understated its recall, so we raised the budget until it answered.

\paragraph{Detector B: a gold-free execution oracle.} For each instance, a model proposes a small
executable check of correct behavior, using only the public task and the visible test---never the gold
patch or the patch under test. We validate the check by running it against the gold patch: a check that the
correct fix fails is unsound and is discarded. We then run each sound check against the candidate. A
candidate that fails a sound check is flagged. All twelve checks were sound. The oracle flags four of the
ten hacks and none of the two correct rewrites.

\begin{table}[t]
\centering
\small
\caption{Two detectors on the same ten certified-resolved reward hacks, with the gold fix withheld from
both. ``Equivalent FP'' is false positives on the two correct rewrites. The judge figures are one
stochastic run.}
\label{tab:detectors}
\begin{tabular}{@{}lccc@{}}
\toprule
Detector & Recall (of 10) & Gold false positives (of 10) & Equivalent FP (of 2) \\
\midrule
A: judge panel (2 models)      & 7 & 1 & 1 \\
B: gold-free execution oracle  & 4 & 0 & 0 \\
\bottomrule
\end{tabular}
\end{table}

\paragraph{The tradeoff.} Table~\ref{tab:detectors} shows the two detectors. The judge reads the patch and
catches more hacks (seven of ten), but it also flags correct code (one gold patch, one correct rewrite).
The execution oracle is blind to the patch and to the gold fix and catches fewer (four of ten), but it
never flagged correct code in this set. The judge trades precision for recall; the oracle trades recall for
precision.

\paragraph{Complementarity.} The two detectors do not catch the same hacks. The judge caught seven; the
oracle caught four; the oracle caught \texttt{django-11211}, which the judge missed. Their union is eight
of ten. Only two hacks, both in astropy, were missed by both. Neither detector alone is enough, and an
ensemble of the two beats either one.

\paragraph{Why the oracle is precise but limited.} The same example explains both halves of the oracle's
behavior. For \texttt{django-11179}, the gold-free check never saw the gold fix or the patch, so it did not
know that the bug depends on the primary key's name. It built a model with a key named \texttt{id}, where
the variant behaves correctly, so its check passed and it missed the hack. The judge caught the same hack
by reading the suspicious \texttt{pk.name == 'id'} condition in the patch. This is why the oracle's recall
is lower: without the gold fix or the patch, a proposed check rarely targets the exact branch where a
certified variant diverges. It is also why the oracle's precision is perfect here: when a spec-derived
check does fire, it fires on real wrong behavior.

\paragraph{Claim boundary.} This is a bounded pilot comparison on ten hacks and two hard controls across
three repositories. The judge result is a single stochastic run. We make no leaderboard, model-superiority,
or state-of-the-art claim.

\section{Limitations}

We state the bounds plainly.
\begin{itemize}
  \item \textbf{Small sample.} A benchmark of twenty-two hacks across eight repositories, with the detector
  comparison run on the first ten. Confidence intervals are wide; this is a pilot, not a distributional
  measurement.
  \item \textbf{Mostly elicited.} The benchmark hacks were produced by a model asked to write a wrong
  variant. We found one unprompted case (above), which shows the class is not only an artifact of
  elicitation, but a single case is an existence result, not a measure of how often a trained agent
  produces such patches on its own.
  \item \textbf{Reference definition of wrong.} ``Wrong'' means the variant differs from the gold reference
  fix on a constructed input, not that an independent oracle judged it wrong.
  \item \textbf{Single-run judge.} The judge numbers come from one run; language-model verdicts vary
  between runs.
  \item \textbf{Coverage-bound oracle.} The gold-free oracle is demonstrated on framework-internal
  functions and is limited by whether a correct check can be proposed from the public specification.
  \item \textbf{One adversary model.} A single frontier model both produced the hacks and proposed the
  checks; we do not measure cross-model generalization.
\end{itemize}

\section{Conclusion}

SWE-bench Verified certifies patches that are wrong. We constructed twenty-two such patches across eight
repositories, verified each by execution, and released them with witnesses. Because the tests that define ``resolved'' are exactly the
tests any test-based check would use, no shipped test catches this class; a check must run an input the
tests do not. Of the two detectors we measured, a language-model judge caught more hacks but flagged
correct code, and a gold-free execution check caught fewer but flagged none. The two are complementary, and
the practical recommendation is an ensemble: a judge to read the patch and a precise execution check to run
it. We found one naturally-occurring certified-yet-wrong patch, produced by a model that was not told to
game the tests; the honest next steps are to grow the sample and the number of repositories, to measure how
often such patches arise unprompted rather than merely showing that they can, and to formalize the
ensemble.

\paragraph{Reproducibility.} Every number in this paper regenerates from a committed script, manifest, or
validated receipt in the Telos repository. The experiments are named in the text (for example, the
construction is iter193 through iter195, the cross-repository expansion is iter199, and the detector
comparison is iter196 and iter197).

\begin{thebibliography}{9}
\small
\bibitem{amodei2016concrete} D.~Amodei, C.~Olah, J.~Steinhardt, P.~Christiano, J.~Schulman, and D.~Man\'e.
Concrete Problems in AI Safety. \emph{arXiv:1606.06565}, 2016.
\bibitem{krakovna2020specification} V.~Krakovna et al. Specification gaming: the flip side of AI ingenuity.
DeepMind Blog, 2020.
\bibitem{skalse2022defining} J.~Skalse, N.~H.~R.~Howe, D.~Krasheninnikov, and D.~Krueger. Defining and
Characterizing Reward Hacking. \emph{Advances in Neural Information Processing Systems (NeurIPS)}, 2022.
\bibitem{jimenez2024swebench} C.~E.~Jimenez, J.~Yang, A.~Wettig, S.~Yao, K.~Pei, O.~Press, and
K.~Narasimhan. SWE-bench: Can Language Models Resolve Real-World GitHub Issues? \emph{Proc.\ ICLR}, 2024.
\emph{arXiv:2310.06770}.
\bibitem{zheng2023judging} L.~Zheng et al. Judging LLM-as-a-Judge with MT-Bench and Chatbot Arena.
\emph{Advances in Neural Information Processing Systems (NeurIPS)}, 2023. \emph{arXiv:2306.05685}.
\bibitem{terekhov2025controltax} M.~Terekhov, Z.~N.~D.~Liu, C.~Gulcehre, and S.~Albanie. Control Tax: The
Price of Keeping AI in Check. \emph{arXiv:2506.05296}, 2025.
\bibitem{deepmind2025fsf} Google DeepMind. Frontier Safety Framework, version 3.0. September 22, 2025.
\bibitem{metr2025malt} METR. MALT: A Dataset of Natural and Prompted Behaviors That Threaten Eval
Integrity. October 14, 2025. (MALT: Manually-reviewed Agentic Labeled Transcripts; 10{,}919 transcripts,
403 tasks, 21 models.)
\bibitem{lambert2024rewardbench} N.~Lambert et al. RewardBench: Evaluating Reward Models for Language
Modeling. \emph{arXiv:2403.13787}, 2024.
\bibitem{chen2018metamorphic} T.~Y.~Chen et al. Metamorphic Testing: A Review of Challenges and
Opportunities. \emph{ACM Computing Surveys}, 51(1), 2018.
\bibitem{claessen2000quickcheck} K.~Claessen and J.~Hughes. QuickCheck: a lightweight tool for random
testing of Haskell programs. \emph{Proc.\ ICFP}, 2000.
\end{thebibliography}

\end{document}
